% GENERATED FILE. Edit canonical lessons, then run npm run generate:books.
% canonical-source-hash: fnv1a64:e18952c21c4391d4
% canonical-lessons: UR-C264-voh-gaya, UR-C264-voh-yahan-tha, UR-R264-where-i-went

\chapter{Where I Went, and Where I Was}
\label{ch:ur-where-i-went-and-where-i-was}
\hlchaptermodality{264}

\begin{chapteropening}
\textbf{By the end of this chapter:} \emph{I can say where someone went and where someone was (voh gayā, voh gaī, maiṅ ghar par thā), and what I used to do (jātā thā).}

Complete UR-R264-where-i-went: say where someone went, where someone was, and what you used to do.
\end{chapteropening}

\section[\texorpdfstring{\ur{وہ} \ur{گیا} (voh gayā)}{وہ گیا (voh gayā)}]{\ur{وہ} \ur{گیا} (voh gayā) — the past, by the one who went}

\label{lesson:UR-C264-voh-gaya}



\begin{quote}
Before the new one: say \textquotedblleft{}pleased to meet you\textquotedblright{}. (\emph{āp se mil kar khushī huī}.) Its last word, \textbf{\ur{ہوئی}} (\emph{huī}), \textquotedblleft{}happened\textquotedblright{}, is a past. This chapter writes the past in Urdu script.
\end{quote}

\subsection*{You'll want to know: the past follows the one who went}
\begin{quote}
\textbf{\ur{وہ} \ur{گیا۔}}
\end{quote}

\begin{quote}
\emph{voh gayā.} — \textquotedblleft{}He went.\textquotedblright{}
\end{quote}

\begin{quote}
\textbf{\ur{وہ} \ur{گئی۔}}
\end{quote}

\begin{quote}
\emph{voh gaī.} — \textquotedblleft{}She went.\textquotedblright{}
\end{quote}

\textbf{\ur{گیا}} (\emph{gayā}) is \textbf{\ur{جانا}} (\emph{jānā}) in the past, with no \textbf{\ur{ہے}} (\emph{hai}) after it. Its ending follows the one who went, as \textbf{\ur{اچھا}} (\emph{achchhā}) follows its noun:

\begin{quote}
\textbf{\ur{وہ} \ur{گھر} \ur{گیا۔}}
\end{quote}

\begin{quote}
\emph{voh ghar gayā.} — \textquotedblleft{}He went home.\textquotedblright{}
\end{quote}

\begin{quote}
\textbf{\ur{سارا} \ur{گھر} \ur{گئی۔}}
\end{quote}

\begin{quote}
\emph{Sārā ghar gaī.} — \textquotedblleft{}Sara went home.\textquotedblright{}
\end{quote}

\begin{quote}
\textbf{\ur{ہم} \ur{بازار} \ur{گئے۔}}
\end{quote}

\begin{quote}
\emph{ham bāzār gae.} — \textquotedblleft{}We went to the market.\textquotedblright{}
\end{quote}

\noindent
\begin{tabularx}{\linewidth}{@{}>{\raggedright\arraybackslash}X>{\raggedright\arraybackslash}X@{}}
\toprule
\textbf{who went} & \textbf{the past} \\
\midrule
a man & \textbf{\ur{گیا}} (\emph{gayā}) \\
a woman & \textbf{\ur{گئی}} (\emph{gaī}) \\
men; \textbf{\ur{ہم}}; \textbf{\ur{آپ}} & \textbf{\ur{گئے}} (\emph{gae}) \\
women & \textbf{\ur{گئیں}} (\emph{gaīṅ}) \\
\bottomrule
\end{tabularx}

\subsection*{You'll want to know: how a verb takes the ending}
Most verbs put the ending on the stem: \textbf{\ur{بیٹھنا}} (\emph{baiṭhnā}) gives \textbf{\ur{بیٹھا}} (\emph{baiṭhā}) and \textbf{\ur{بیٹھی}} (\emph{baiṭhī}).

\begin{quote}
\textbf{\ur{وہ} \ur{بیٹھی۔}}
\end{quote}

\begin{quote}
\emph{voh baiṭhī.} — \textquotedblleft{}She sat down.\textquotedblright{}
\end{quote}

A stem that ends in a vowel adds \textbf{-\ur{یا}} for a man and \textbf{-\ur{ئی}} for a woman:

\begin{quote}
\textbf{\ur{وہ} \ur{کل} \ur{آیا۔}}
\end{quote}

\begin{quote}
\emph{voh kal āyā.} — \textquotedblleft{}He came yesterday.\textquotedblright{}
\end{quote}

\begin{quote}
\textbf{\ur{سارا} \ur{کل} \ur{آئی۔}}
\end{quote}

\begin{quote}
\emph{Sārā kal āī.} — \textquotedblleft{}Sara came yesterday.\textquotedblright{}
\end{quote}

\textbf{\ur{سونا}} (\emph{sonā}) gives \textbf{\ur{سویا}} (\emph{soyā}), and \textbf{\ur{اٹھنا}} (\emph{uṭhnā}) gives \textbf{\ur{اٹھا}} (\emph{uṭhā}). \textbf{\ur{ہونا}} (\emph{honā}) gives \textbf{\ur{ہوا}} (\emph{huā}) and \textbf{\ur{ہوئی}} (\emph{huī}): \textbf{\ur{کیا} \ur{ہوا؟}} (\emph{kyā huā?}), \textquotedblleft{}what happened?\textquotedblright{} A bare stem with \textbf{\ur{گیا}} says a change is complete: \textbf{\ur{وہ} \ur{سو} \ur{گیا}} (\emph{voh so gayā}), \textquotedblleft{}he fell asleep\textquotedblright{}. \textquotedblleft{}Did not\textquotedblright{} puts \textbf{\ur{نہیں}} (\emph{nahīṅ}) before the verb:

\begin{quote}
\textbf{\ur{میں} \ur{نہیں} \ur{گیا۔}}
\end{quote}

\begin{quote}
\emph{maiṅ nahīṅ gayā.} — \textquotedblleft{}I didn't go.\textquotedblright{}
\end{quote}

\subsection*{Your turn}
Say these aloud:

\begin{itemize}
\raggedright
  \item \textquotedblleft{}He came home\textquotedblright{} — \emph{voh ghar āyā}
  \item \textquotedblleft{}She came home\textquotedblright{} — \emph{voh ghar āī}
  \item \textquotedblleft{}We came home\textquotedblright{} — \emph{ham ghar āe}
\end{itemize}

\begin{tcolorbox}[breakable,colback=teal!4,colframe=teal!35!black,title={Before you move on}]
\textquotedblleft{}She went\textquotedblright{}? (\textbf{\ur{وہ} \ur{گئی}}.) \textquotedblleft{}He sat down\textquotedblright{}? (\textbf{\ur{وہ} \ur{بیٹھا}}.) \textquotedblleft{}What happened?\textquotedblright{} (\textbf{\ur{کیا} \ur{ہوا؟}}) Who decides the ending? (The one who went.)
\end{tcolorbox}
\section[\texorpdfstring{\ur{وہ} \ur{یہاں} \ur{تھا} (voh yahāṅ thā)}{وہ یہاں تھا (voh yahāṅ thā)}]{\ur{وہ} \ur{یہاں} \ur{تھا} (voh yahāṅ thā) — was, were, used to}

\label{lesson:UR-C264-voh-yahan-tha}



\begin{quote}
Say \textquotedblleft{}he went\textquotedblright{} and \textquotedblleft{}she came\textquotedblright{}. (\emph{voh gayā}, \emph{voh āī}.) Then \textquotedblleft{}he is here\textquotedblright{}. (\emph{voh yahāṅ hai}.) Now put \textbf{\ur{ہے}} (\emph{hai}) into the past.
\end{quote}

\subsection*{You'll want to know: \ur{تھا}, \ur{تھی}, \ur{تھے}}
\begin{quote}
\textbf{\ur{وہ} \ur{یہاں} \ur{تھا۔}}
\end{quote}

\begin{quote}
\emph{voh yahāṅ thā.} — \textquotedblleft{}He was here.\textquotedblright{}
\end{quote}

\begin{quote}
\textbf{\ur{وہ} \ur{یہاں} \ur{تھی۔}}
\end{quote}

\begin{quote}
\emph{voh yahāṅ thī.} — \textquotedblleft{}She was here.\textquotedblright{}
\end{quote}

\textbf{\ur{تھا}} (\emph{thā}) is \textbf{\ur{ہے}} in the past. Like \textbf{\ur{گیا}}, it follows the person:

\begin{quote}
\textbf{\ur{میں} \ur{گھر} \ur{پر} \ur{تھا۔}}
\end{quote}

\begin{quote}
\emph{maiṅ ghar par thā.} — \textquotedblleft{}I was at home.\textquotedblright{}
\end{quote}

\begin{quote}
\textbf{\ur{سارا} \ur{بیمار} \ur{تھی۔}}
\end{quote}

\begin{quote}
\emph{Sārā bīmār thī.} — \textquotedblleft{}Sara was ill.\textquotedblright{}
\end{quote}

\begin{quote}
\textbf{\ur{ہم} \ur{اسکول} \ur{میں} \ur{تھے۔}}
\end{quote}

\begin{quote}
\emph{ham iskūl meṅ the.} — \textquotedblleft{}We were at school.\textquotedblright{}
\end{quote}

\noindent
\begin{tabularx}{\linewidth}{@{}>{\raggedright\arraybackslash}X>{\raggedright\arraybackslash}X@{}}
\toprule
\textbf{who} & \textbf{was} \\
\midrule
a man & \textbf{\ur{تھا}} (\emph{thā}) \\
a woman & \textbf{\ur{تھی}} (\emph{thī}) \\
men; \textbf{\ur{ہم}}; \textbf{\ur{آپ}} & \textbf{\ur{تھے}} (\emph{the}) \\
women & \textbf{\ur{تھیں}} (\emph{thīṅ}) \\
\bottomrule
\end{tabularx}

\begin{quote}
\textbf{\ur{کل} \ur{آپ} \ur{کہاں} \ur{تھے؟}}
\end{quote}

\begin{quote}
\emph{kal āp kahāṅ the?} — \textquotedblleft{}Where were you yesterday?\textquotedblright{}
\end{quote}

\subsection*{You'll want to know: what I used to do}
Keep the present \textbf{\ur{جاتا}} (\emph{jātā}) and put \textbf{\ur{تھا}} where \textbf{\ur{ہے}} stood. That is a habit in the past, \textquotedblleft{}used to\textquotedblright{}:

\begin{quote}
\textbf{\ur{میں} \ur{اسکول} \ur{جاتا} \ur{تھا۔}}
\end{quote}

\begin{quote}
\emph{maiṅ iskūl jātā thā.} — \textquotedblleft{}I used to go to school.\textquotedblright{}
\end{quote}

\begin{quote}
\textbf{\ur{وہ} \ur{ہر} \ur{روز} \ur{آتی} \ur{تھی۔}}
\end{quote}

\begin{quote}
\emph{voh har roz ātī thī.} — \textquotedblleft{}She used to come every day.\textquotedblright{}
\end{quote}

Both halves follow the person: \textbf{\ur{جاتی} \ur{تھی}} (\emph{jātī thī}) for a woman, \textbf{\ur{جاتے} \ur{تھے}} (\emph{jāte the}) for men.

\subsection*{Your turn}
Say these aloud:

\begin{itemize}
\raggedright
  \item \textquotedblleft{}I was at the office\textquotedblright{} — \emph{maiṅ daftar meṅ thā}
  \item \textquotedblleft{}She was at the office\textquotedblright{} — \emph{voh daftar meṅ thī}
  \item \textquotedblleft{}We were at the office\textquotedblright{} — \emph{ham daftar meṅ the}
\end{itemize}

\begin{tcolorbox}[breakable,colback=teal!4,colframe=teal!35!black,title={Before you move on}]
\textquotedblleft{}She was ill\textquotedblright{}? (\textbf{\ur{وہ} \ur{بیمار} \ur{تھی}}.) \textquotedblleft{}We were at home\textquotedblright{}? (\textbf{\ur{ہم} \ur{گھر} \ur{پر} \ur{تھے}}.) \textquotedblleft{}I used to go\textquotedblright{}? (\textbf{\ur{میں} \ur{جاتا} \ur{تھا}}, or \textbf{\ur{جاتی} \ur{تھی}}.) Which word decides \textbf{\ur{تھا}} or \textbf{\ur{تھی}}? (The person.)
\end{tcolorbox}
\section[(review)]{Review — where I went, and where I was}

\label{lesson:UR-R264-where-i-went}



\begin{quote}
Say \textquotedblleft{}he went\textquotedblright{} and \textquotedblleft{}she was here\textquotedblright{}. (\emph{voh gayā}, \emph{voh yahāṅ thī}.)
\end{quote}

\subsection*{Your turn}
Say these aloud:

\begin{itemize}
\raggedright
  \item Yesterday I went to the market. Sara went home.
  \item He came, and then he sat down.
  \item I didn't go. I was ill.
  \item Where were you yesterday? We were at home.
  \item I used to go to school every day.
\end{itemize}

\begin{tcolorbox}[breakable,colback=teal!4,colframe=teal!35!black,title={Before you move on}]
\textbf{\ur{وہ} \ur{گئی}} or \textbf{\ur{وہ} \ur{گیا}}, for Sara? (\textbf{\ur{گئی}}.) Now the other way round: what does \textbf{\ur{تھے}} tell you about the people? (Men, or \textbf{\ur{ہم}}, or \textbf{\ur{آپ}}.) \textquotedblleft{}I used to come\textquotedblright{}? (\textbf{\ur{میں} \ur{آتا} \ur{تھا}}.) In every one, who decides the ending? (The one who went, or who was.)
\end{tcolorbox}
